% Owner: counterfactual_algebra. New attachment only; no frozen edition edited.
% Sources retained: SP.1--8, SP.25--31; TP.1--8, TP.23--29;
% CFA.8--14 and CFA.33--35. New localization calculations: MLT.
\section{The original localization triangle with its marked derived fibre}
\label{sec:MLT}

\paragraph{Original objects and conventions.}
Put \(A=\mathbb C[u]\), \(A^\times_u=\mathbb C[u,u^{-1}]\),
and \(\kappa=A/(u)=\mathbb C\). The superscript in
\(A^\times_u\) denotes this specified localization, not a group
of units. Fix the original monic polynomial \(h(s)\) of degree
\(d\ge1\) and an ordered tensor length \(k\ge1\). The argument
first treats this entire polynomial, with every coefficient retained.
For the arithmetic marking, \(h\) is a nonempty full packet divisor
of \(g=2\xi\): every retained zero has its complete order.
This does not assert that any zero is off the critical line.
The empty packet has zero marked coefficient module and hence zero
lattice, localized module, quotient, and derived fibre in the
formulas for those modules below.
Let
\[
 \begin{gathered}
 R_k=A[s_1,\ldots,s_k],\qquad
 \Phi_h(s)=\int_0^s h(z)\,dz,\qquad
 \Phi_k=\sum_{i=1}^k\Phi_h(s_i),\\
 K_k^j=R_k\otimes_{\mathbb C}\bigwedge\nolimits^j
       \langle ds_1,\ldots,ds_k\rangle,\qquad
 d_K=u\,d+d\Phi_k\wedge
       =\sum_{i=1}^k ds_i\wedge L_i,\qquad
 L_i=u\partial_{s_i}+h(s_i).
 \end{gathered}                                             \tag{MLT.1}
\]
All variables \(s_i\), the primitive's base point \(0\), and the
order of the top form \(ds_1\wedge\cdots\wedge ds_k\) are fixed.
The differential of a tensor is
\(\sum_i(-1)^{\sum_{j<i}|x_j|}
x_1\otimes\cdots\otimes d x_i\otimes\cdots\otimes x_k\).
The operators \(L_i,L_j\) commute for \(i\ne j\), because
each differentiates and multiplies only its own variable.
Their two wedge orders have opposite signs. Consequently
\(d_K^2=0\). All complexes use cohomological degrees, and
\((C[1])^j=C^{j+1}\), with differential \(-d_C\).

\paragraph{The original monic remainder and an explicit retraction.}
For \(Q\ne0\) in \(A[s]\), the leading \(s\)-degree of
\((u\partial_s+h)Q\) is \(\deg_sQ+d\), with exactly the
leading coefficient of \(Q\). The derivative has smaller
degree, and the leading coefficient of \(h\) is \(1\).
For any \(f\in A[s]\), subtract that leading coefficient
times \((u\partial_s+h)s^{\deg_s f-d}\) while its degree
is at least \(d\). Each subtraction strictly decreases
the degree, so the procedure terminates and gives
\[
 f=LQ+\operatorname{Rem}_u(f),\qquad
 \deg_s\operatorname{Rem}_u(f)<d,\qquad
 A[s]=L(A[s])\oplus A[s]_{<d}.                            \tag{MLT.2}
\]
The representation is unique: a nonzero \(LQ\) cannot
have degree less than \(d\). Uniqueness also proves
that the quotient \(Q\) and remainder are \(A\)-linear
functions of \(f\).
For the one-factor complex \(K_1=[A[s]\xrightarrow{L}A[s]ds]\),
let \(V_1=A[s]_{<d}ds\), concentrated in degree one.
Define \(i_1:V_1\to K_1\) by inclusion,
\(p_1:K_1\to V_1\) by zero in degree zero and remainder
in degree one, and \(H_1:K_1^1\to K_1^0\) by
\(H_1(fds)=Q\) from (MLT.2), zero in all other degrees.
Then
\[
 p_1i_1=1,\qquad
 d_KH_1+H_1d_K=1-i_1p_1.                                 \tag{MLT.3}
\]
In degree zero this is \(H_1L=1\); in degree one it
is \(LH_1=1-i_1p_1\), exactly (MLT.2).

Set \(q_1=i_1p_1\). Tensor these maps in the given order.
The projection \(p_k=p_1^{\otimes k}\) is zero below
degree \(k\), and in degree \(k\) takes successive
differential remainders in the separate variables.
The inclusion \(i_k=i_1^{\otimes k}\) inserts those
remainders with the original top form. An explicit
tensor homotopy is
\[
 H_k=\sum_{i=1}^k q_1^{\otimes(i-1)}\otimes H_1
                        \otimes 1^{\otimes(k-i)},\qquad
 d_KH_k+H_kd_K=1-i_kp_k.                                 \tag{MLT.4}
\]
On a homogeneous tensor the \(i\)-th summand has sign
\((-1)^{\sum_{j<i}|x_j|}\), because \(H_1\) has degree
\(-1\). Terms differentiating a different position
cancel in \(d_KH_k+H_kd_K\); the remaining terms are
\(\sum_i q_1^{\otimes(i-1)}\otimes(1-q_1)
\otimes1^{\otimes(k-i)}=1-q_1^{\otimes k}\).
This telescoping calculation proves (MLT.4).

Let \(E_h=\mathbb C[s]/(h)\), and index the original
ordered monomial basis by \(\mathbf b\in\{0,\ldots,d-1\}^k\).
Write \(e_{\mathbf b}=[s_1^{b_1}\cdots s_k^{b_k}
ds_1\wedge\cdots\wedge ds_k]\). The result is the
explicit homotopy equivalence
\[
 \begin{gathered}
 K_k\simeq L[-k],\qquad
 L=\bigoplus_{\mathbf b}A e_{\mathbf b},\qquad
 H^j(K_k)=
 \begin{cases}L&j=k,\\0&j\ne k,\end{cases}\\
 L/uL=
 \frac{\mathbb C[s_1,\ldots,s_k]}
      {(h(s_1),\ldots,h(s_k))}
       ds_1\wedge\cdots\wedge ds_k
 =E_h^{\otimes k}ds_1\wedge\cdots\wedge ds_k.              \tag{MLT.5}
 \end{gathered}
\]
The quotient identification follows by successive ordinary
monic division at \(u=0\). Its basis is precisely the
same ordered monomial list, so it preserves every local
nilpotent coordinate of the original polynomial.
All formulas (MLT.2)--(MLT.4) specialize literally to
the original marked complex at \(u=0\).

\paragraph{The one localization sequence at complex level.}
Localize only the original parameter \(u\) and define
\[
 \begin{gathered}
 K_k^{\mathrm{loc}}=A^\times_u\otimes_A K_k,\qquad
 C_k=K_k^{\mathrm{loc}}/K_k,\qquad
 M=A^\times_u\otimes_A L,\qquad Q=M/L,\\
 0\longrightarrow K_k\xrightarrow{\iota}K_k^{\mathrm{loc}}
       \xrightarrow{\pi}C_k\longrightarrow0 .
 \end{gathered}                                           \tag{MLT.6}
\]
Each term of \(K_k\) is a free \(A\)-module, so its
localization map is injective: if \(u^r x=0\) for a
polynomial coefficient vector, then \(x=0\). This
proves the degreewise short exact sequence in (MLT.6).
The differential is \(A\)-linear, hence preserves
the original lattice and defines both the localized
and quotient differentials. The maps \(i_k,p_k,H_k\)
are also \(A\)-linear and preserve the original
lattices. They therefore localize and descend to
the quotient, with exactly the same homotopy equations.
In particular
\[
 K_k^{\mathrm{loc}}\simeq M[-k],\qquad
 C_k\simeq Q[-k],\qquad
 H^j(C_k)=
 \begin{cases}Q&j=k,\\0&j\ne k.\end{cases}                 \tag{MLT.7}
\]
Thus the full cohomology long exact sequence of
(MLT.6) has sole nonzero part
\[
                 0\longrightarrow L\longrightarrow M
                    \longrightarrow Q\longrightarrow0. \tag{MLT.8}
\]
This is an actual quotient of the original complexes,
not removal of their marked cohomology by a new
support convention.

For precision about the triangle, use
\(\operatorname{Cone}(\iota)^j
=(K_k^{\mathrm{loc}})^j\oplus K_k^{j+1}\) with
\[
 d_{\mathrm{cone}}(b,a)=(d b+\iota(a),-d a),\qquad
 (b,a)\longmapsto[b]\in C_k.                             \tag{MLT.9}
\]
The latter map is a quasi-isomorphism. Its kernel is
identified with pairs \((c,a)\in K_k^j\oplus K_k^{j+1}\)
and differential \((dc+a,-da)\); the degree-minus-one
homotopy \((c,a)\mapsto(0,c)\) has
\(dH+Hd=1\). The quotient therefore yields the
distinguished triangle represented by
\[
           K_k\longrightarrow K_k^{\mathrm{loc}}
                \longrightarrow C_k\longrightarrow K_k[1].
                                                               \tag{MLT.10}
\]
With the specific cone projection \((b,a)\mapsto a\)
as last arrow, its connecting map is the negative of
the raw short-exact-sequence boundary obtained by
lifting a cycle and differentiating it: a lift \(b\)
has cone cycle \((b,-db)\). This sign convention is
kept separate from the positive Tor identification below.

\paragraph{The complete quotient module and its derived marked fibre.}
Every class in \(Q\) has a unique finite principal-part
expansion
\[
 Q=\bigoplus_{r\ge1}u^{-r}E_h^{\otimes k}
                 ds_1\wedge\cdots\wedge ds_k,\qquad
 u[u^{-r}v]=
 \begin{cases}[u^{-r+1}v]&r\ge2,\\0&r=1.\end{cases}        \tag{MLT.11}
\]
Here \(v\) is a constant coefficient vector in the
fixed ordered basis; the direct sum means finite
support in negative powers. Existence and uniqueness
are those of a Laurent polynomial after deleting
its nonnegative powers. In particular multiplication
by \(u\) is surjective on \(Q\), and
\[
 Q/uQ=0,\qquad Q[u]=\{[u^{-1}v]:v\in L\},\qquad
 \delta:Q[u]\xrightarrow{\ \sim\ }L/uL,\quad
       [u^{-1}v]\longmapsto v\bmod uL.                   \tag{MLT.12}
\]
To check the last formula without choosing a
principal-part representative, a class \([m]\in Q[u]\)
has \(um=v\in L\), so \(m=u^{-1}v\).
Changing \(m\) by \(\ell\in L\) changes \(v\) by
\(u\ell\), exactly the stated ambiguity.
Conversely \(v\bmod uL\mapsto[u^{-1}v]\) is well-defined
and inverse. This proves both injectivity and
surjectivity of the displayed map.

We calculate derived specialization with the free
resolution
\[
 F=[A\epsilon_{-1}\xrightarrow{\ u\ }A\epsilon_0],
 \qquad |\epsilon_{-1}|=-1,\quad|\epsilon_0|=0,
 \qquad F\longrightarrow\kappa,                          \tag{MLT.13}
\]
placed as the \emph{left} tensor factor throughout.
For a complex \(X\), the total differential is
\(d(f\otimes x)=d_Ff\otimes x+(-1)^{|f|}f\otimes d_Xx\).
The resolution is exact in degree \(-1\), because
multiplication by \(u\) is injective on \(A\), and
has cokernel \(\kappa\) in degree zero. Tensoring
this two-term free resolution computes derived
specialization: it preserves quasi-isomorphisms,
since the tensor of a bounded complex of free
modules with an acyclic complex is acyclic, as
follows by its finite filtration with successive
quotients copies of that acyclic complex.
For a module \(V\), this calculation gives
\(H^{-1}(F\otimes_A V)=\ker(u:V\to V)\) and
\(H^0(F\otimes_A V)=V/uV\). Applying it to
(MLT.12) proves
\[
 \operatorname{Tor}_1^A(Q,\kappa)
   \xrightarrow[\ [u^{-1}v]\mapsto v\bmod uL\ ]{\ \delta\ }
   L/uL
   =E_h^{\otimes k}ds_1\wedge\cdots\wedge ds_k,\qquad
 \operatorname{Tor}_j^A(Q,\kappa)=0\quad(j\ne1).
                                                               \tag{MLT.14}
\]
Here \(\operatorname{Tor}_0=Q\otimes_A\kappa=0\)
is included in the vanishing statement.
Thus the claimed isomorphism is an actual boundary
on the entire packet and includes every basis vector
and every linear combination, not only its dimension.

For \(M\), multiplication by \(u\) is invertible.
The complex \(F\otimes_A M\) contracts by sending
\(\epsilon_0\otimes m\) to
\(\epsilon_{-1}\otimes u^{-1}m\), and sending
\(\epsilon_{-1}\otimes m\) to zero. The same formula
contracts \(F\otimes_A K_k^{\mathrm{loc}}\);
the terms involving its internal differential
cancel by the left tensor sign convention.
For \(L\), multiplication by \(u\) is injective,
so its derived fibre is its ordinary fibre in
degree zero. Combining these explicit calculations
with the retained retractions gives
\[
 \begin{aligned}
 F\otimes_A K_k&\simeq (L/uL)[-k],\\
 F\otimes_A K_k^{\mathrm{loc}}&\simeq0,\\
 F\otimes_A C_k&\simeq (L/uL)[-(k-1)].
 \end{aligned}                                           \tag{MLT.15}
\]
The last complex therefore has its sole cohomology
in degree \(k-1\). The quotient's ordinary
specialization is zero, but its derived
specialization retains exactly the original marked
cohomology one degree lower.

The positive connecting formula can be seen
directly in the original complex. For \(v\in L\),
the element
\(\epsilon_{-1}\otimes[u^{-1}i_k(v)]\)
is a degree-\(k-1\) cycle in \(F\otimes_A C_k\).
Lift it to
\(\epsilon_{-1}\otimes u^{-1}i_k(v)\)
in \(F\otimes_A K_k^{\mathrm{loc}}\).
Its differential is
\(\epsilon_0\otimes i_k(v)\), since \(i_k(v)\)
is a top-degree cycle. Hence the raw long exact
sequence boundary of the specialized short exact
sequence is exactly
\[
 [\epsilon_{-1}\otimes u^{-1}i_k(v)]
       \longmapsto[\epsilon_0\otimes i_k(v)]
       \longleftrightarrow v\bmod uL.                   \tag{MLT.16}
\]
There is no parity sign in (MLT.16), because \(F\)
is the left factor. The cone arrow of (MLT.9)
induces the negative of this raw boundary, as
already specified. Both descriptions retain
the original coordinate and its stated sign.

\paragraph{All pole orders, coefficient maps, and strict complex lifts.}
The first-order calculation also retains each
finite higher thickening. For \(n\ge1\),
\[
 Q[u^n]=u^{-n}L/L,\qquad
 \delta_n:Q[u^n]\xrightarrow{\ \sim\ }L/u^nL,\quad
       [u^{-n}v]\mapsto v\bmod u^nL,\qquad Q/u^nQ=0.
                                                               \tag{MLT.17}
\]
The proof of (MLT.12) applies verbatim with \(u^n\):
a class killed by \(u^n\) has the prescribed
representative, and changing a lift changes its
numerator by \(u^nL\). Surjectivity of \(u^n\)
follows by increasing every pole order by \(n\).
For \(n\le N\), inclusion \(Q[u^n]\hookrightarrow Q[u^N]\)
corresponds under these maps to
\[
 L/u^nL\longrightarrow L/u^NL,\qquad
       v\bmod u^nL\longmapsto u^{N-n}v\bmod u^NL.
                                                               \tag{MLT.18}
\]
Indeed \(u^{-n}v=u^{-N}u^{N-n}v\); this also
proves well-definedness of the displayed map.

Let \(T(u):L\to L'\) be any of the specified
coefficient-linear maps in the original remainder
bases, extended \(A\)-linearly, or more generally
any matrix with entries in \(A\). It has the
unique localized extension
\(u^{-r}v\mapsto u^{-r}T(u)v\), which preserves
the target lattice and induces a map \(Q\to Q'\).
These maps give a commuting map between the
short exact sequences (MLT.8), and
\[
 \delta' \bigl([u^{-1}T(u)v]\bigr)
       =T(0)(v\bmod uL),\qquad
 \delta_N'\bigl([u^{-N}T(u)v]\bigr)
       =T(u)v\bmod u^NL'.                               \tag{MLT.19}
\]
The first equality follows by reducing \(T(u)\)
modulo \(u\), and the second by reducing it
modulo \(u^N\). Thus all these diagrams are
natural for the actual matrices, with their full
coefficients and source and target labels.

For such a map \(T\), the original retractions
give the explicit cochain map
\[
 \widetilde T=i_k'Tp_k:K_k\longrightarrow K_k',
 \qquad
 \widetilde T\,\widetilde S=\widetilde{TS}.
                                                               \tag{MLT.20}
\]
The projection and inclusion are cochain maps,
so this is a cochain map; \(p_k'i_k'=1\)
proves the composition formula when the targets
are matched. This lift is linear in \(T\).
Its value at the identity is \(i_kp_k\), with
the explicit homotopy to \(1\) in (MLT.4).
For an invertible coefficient matrix \(U\)
on a fixed lattice there is also a literal
cochain automorphism, retaining the contractible
summand:
\[
 \widehat U=i_kUp_k+(1-i_kp_k),\qquad
 \widehat U^{-1}=i_kU^{-1}p_k+(1-i_kp_k).                 \tag{MLT.21}
\]
Both terms commute with the differential.
The projectors \(i_kp_k\) and \(1-i_kp_k\)
are complementary. The two displayed summands
act on their respective images and their
cross products vanish; \(p_ki_k=1\) then
proves the inverse formula. These maps preserve the
original polynomial lattice and therefore
act on the entire localization sequence and
on its left-resolution specialization.

There is an exact comparison with any original
cochain arrow \(f:K_k\to K_k'\) between these
presentations. Put \(T=p_k'f i_k\) and
\[
 \mathcal H=H_k'f+(i_k'p_k')fH_k.
 \qquad
 f-i_k'Tp_k=d\mathcal H+\mathcal H d.                    \tag{MLT.22}
\]
Expand the right side using (MLT.4) and the
cochain property of \(f\). The result is
\((1-i_k'p_k')f+(i_k'p_k')f(1-i_kp_k)
=f-(i_k'p_k')f(i_kp_k)\), the left side.
Consequently the matrix lift records the
exact induced map and a proved homotopy;
it does not replace an original representative
map by an unproved equality.

\paragraph{The full arithmetic unit and the original theta marking.}
For a full packet \(h\), retain
\[
 v_h=g/h,\qquad \upsilon_h=j_hv_h\in E_h^\times,\qquad
 \varepsilon_h=\upsilon_h^{-1},\qquad
 U_h=M_{\upsilon_h}\text{ in }(1,s,\ldots,s^{d-1}),
 \qquad U_k=U_h^{\otimes k}.                              \tag{MLT.23}
\]
Here \(j_h\) means all the Taylor jets to the
complete orders specified by \(h\).
At each retained centre \(\rho\) of order \(m_\rho\),
the function \(v_h\) has nonzero constant Taylor
coefficient. Writing
\(v_h(\rho+z)=\sum_{j=0}^{m_\rho-1}a_{\rho,j}z^j\)
modulo \(z^{m_\rho}\), the full inverse is given by
\[
 b_{\rho,0}=a_{\rho,0}^{-1},\qquad
 b_{\rho,r}=-a_{\rho,0}^{-1}
              \sum_{j=1}^r a_{\rho,j}b_{\rho,r-j}
                  \quad(1\le r<m_\rho).                 \tag{MLT.24}
\]
Multiplication of the two finite series gives
constant coefficient one and each other
coefficient zero, proving the inverse in that
local quotient. The factors \((s-\rho)^{m_\rho}\)
are pairwise coprime. Bezout identities from
polynomial division give complementary
idempotents and reconstruct a unique class
in \(E_h\) from all its local classes. Thus
(MLT.24) constructs \(\varepsilon_h\) and
proves invertibility of the full matrix \(U_h\),
without discarding any nilpotent coefficient.
The tensor inverse is \((U_h^{-1})^{\otimes k}\).
These constant coefficient matrices act on
\(L,M,Q\) by (MLT.19), and on the complexes
by (MLT.21).

The source maps are the original ones retained
in CFA.9--14 and SP.27:
\[
 \begin{gathered}
 D=-x\partial_x,\quad
 \Theta\phi(x)=2\sum_{n\ge1}\phi(nx),\quad
 \phi_*=(4\pi^2x^4-6\pi x^2)e^{-\pi x^2},\\
 h(D)F_h=\Theta\phi_*,\qquad
 \mathcal MF_h=v_h,\qquad
 \alpha_h(P)=P(D)\phi_*,\quad\mathcal T_h(P)=P(D)F_h,\\
 \Theta\alpha_h(P)=\mathcal T_h(hP),\qquad
 J_h\mathcal T_h(P)=U_h[P]_s,\qquad
 s_h(q)=\mathcal T_h(r_h(\varepsilon_h q)),\qquad
 \sigma_h=[s_h]\in\mathscr B/\Theta V .
 \end{gathered}                                             \tag{MLT.25}
\]
Here \(r_h\) is the original monic remainder
representative. The Mellin convention is
\(\mathcal MF(s)=\int_0^\infty F(x)x^s\,dx/x\);
the spaces are the original even Schwartz
\(V\), with \(\phi(0)=\int\phi=0\), and the
original \(\mathscr B\) with all seminorms
\(\sup_{x>0}x^b|D^rF(x)|\) finite for integral
\(b\) and nonnegative \(r\).
The original Euler inverse constructs \(F_h\);
its uniqueness and the source identities are
those proved in the retained source sections.
No different inverse or source space is
introduced here. The new maps below follow
directly from these actual identities.

Indeed integration by parts, whose endpoints
vanish by those seminorms, gives
\(\mathcal M(P(D)F_h)=P(s)v_h(s)\).
Its complete jet is \(U_h[P]_s\).
Also \(h(D)F_h=\Theta\phi_*\) and the
commutation of \(D\) with \(\Theta\) give
the polynomial boundary identity in (MLT.25).
It follows that \([\mathcal T_h]\) factors
through \(E_h\) into the original theta
cohomology. Its corrected section has
\(J_hs_h(q)=\upsilon_h\varepsilon_hq=q\),
so it is injective after passage to that
cohomology, where \(J_h\) vanishes on
the original theta boundaries.
For ordered tensors the actual source
map on the marked coefficient column \(v\) is
\[
 \tau_k(v)=\sigma_h^{\otimes k}U_kv,\qquad
 q=U_kv,\qquad
 [u^{-1}v]\xrightarrow{\delta}v
       \xrightarrow{U_k}q
       \xrightarrow{\sigma_h^{\otimes k}}\tau_k(v).
                                                               \tag{MLT.26}
\]
The tensor source square has precisely the
ordered signs of (MLT.1). A tensor of the
injective sections is injective: tensor
their linear left inverses \(J_h\) on
their images. Thus every nonzero original
marked vector gives a nonzero Tor class,
and (MLT.26) recovers its exact original
theta class. In the corrected coordinate
\(q\), the representative is
\([u^{-1}U_k^{-1}q]\), with the displayed
full inverse unit retained.

The same maps include the original divisor
arrows. For full \(h\mid H\), write \(b=H/h\).
Old and new factors are disjoint, so their
coprimality defines restriction \(\pi:E_H\to E_h\)
and arithmetic zero insertion \(\iota:E_h\to E_H\)
by their complete local jets. The raw numerator
map is
\[
 \mu_b([P]_h)=[bP]_H
       =\iota M_{[b]_h}([P]_h),\qquad
 \iota=\mu_bM_{[b]_h^{-1}}.                               \tag{MLT.27}
\]
On every old factor these formulas are
multiplication by the full unit \(b\);
on a new factor \(b\) vanishes to its
complete order. This proves both equations
and their well-definedness.
Moreover \(v_h=bv_H\), whence
\(\pi\varepsilon_H=[b]_h\varepsilon_h\).
The polynomial \(b\,r_h(\varepsilon_hq)\)
has degree below \(\deg H\) and complete
jet \(\varepsilon_H\iota q\). Uniqueness
of the monic remainder therefore gives
\[
 r_H(\varepsilon_H\iota q)=b\,r_h(\varepsilon_hq),
 \qquad s_H\iota=s_h,\qquad \sigma_H\iota=\sigma_h .
                                                               \tag{MLT.28}
\]
For the source equality use the original
identity \(b(D)F_H=F_h\), proved by applying
the injective operator \(h(D)\) on \(\mathscr B\).
These matrices, their ordered tensor powers,
and their full unit factors extend through
(MLT.19)--(MLT.22). Thus the Tor diagrams
commute with the original restriction,
zero insertion, and raw numerator arrows
with their distinct types.

\paragraph{Representative multiplication and its entire defect.}
Let \(\nu_h=r_h(\upsilon_h)\) be the original
degree-less-than-\(d\) representative of the
full unit jet, and \(V_{\mathrm{pol}}(s)
=\prod_i\nu_h(s_i)\). For every polynomial
\(f\), the actual twisted differential satisfies
\[
 V_{\mathrm{pol}}L_if
    =L_i(V_{\mathrm{pol}}f)
       -u(\partial_{s_i}V_{\mathrm{pol}})f.               \tag{MLT.29}
\]
This is the Leibniz rule in its original
variable; both multiplication terms agree
and the displayed derivative is the entire
remaining term. Thus the constant matrix
\(U_k\) above is a specified coefficient
extension, not an assertion that multiplying
twisted representatives is a cochain
endomorphism.
For \(0\le b<d\), perform the original
ordinary monic division
\[
 \nu_h(s)s^b=h(s)q_b(s)+r_b(s),\qquad
 \deg r_b<d,\quad \deg q_b\le d-2.
                                                               \tag{MLT.30}
\]
The zero quotient is allowed, including
\(d=1\). Since \(q_b'\) already has degree
less than \(d\), subtracting \(Lq_b\)
shows that its exact differential remainder
is \(r_b-uq_b'\).
Let \(H_h\) have column \(b\) the original
coefficient column of \(q_b'\).
Successive separate-variable reductions give
\[
 \operatorname{Rem}_u(V_{\mathrm{pol}}\cdot)
          =(U_h-uH_h)^{\otimes k}
          \quad\text{on the ordered remainder frame}.   \tag{MLT.31}
\]
For one factor this is precisely the previous
subtraction; independent variables and the
unchanged top-form order prove the tensor
formula. This exact correction is retained
alongside the coefficient automorphism in
(MLT.21). They agree at \(u=0\) and are
related by (MLT.29)--(MLT.31) at general
\(u\); no derivative term has been omitted.

\paragraph{Every independent coefficient face and its source label.}
Fix a finite index set \(I\) of original
coefficient positions, retaining each position's
original source-leg label. For every
\(S\subseteq I\), form the finite direct sums
of (MLT.6) and (MLT.8) over the active positions:
\[
 \begin{gathered}
 0\to K_k[S]\to K_k^{\mathrm{loc}}[S]\to C_k[S]\to0,\\
 0\to L[S]\to M[S]\to Q[S]\to0,\qquad
 \delta[S]:Q[S][u]\xrightarrow{\sim}L[S]/uL[S].
 \end{gathered}                                           \tag{MLT.32}
\]
An empty active set gives zero modules.
The underlying active face consists of all
tuples indexed by \(S\), with every coordinate
amplitude retained, including zero. An
inactive position is the tagged symbol
\(\tau_i\); an active zero is the supported
zero \(e_i\). These tags are distinct.
Thus the actual face carrier is the
disjoint union over \(S\) of the indicated
active tuples with the inactive tags
adjoined in their original positions.
At the marked fibre it is exactly the
independent split carrier
\(\prod_{i\in I}G(E_h^{\otimes k})\),
with each specified source label attached.
This construction does not give \(Q\)
a quotient ring structure: \(L\) is an
additive submodule of \(M\), not a ring
ideal, and all arrows involving \(Q\)
in (MLT.32) are module and labelled-face
maps of the stated types.

For \(S\subseteq T\), the face insertion
fills every newly active position of
\(T\setminus S\) with its supported zero.
On active coefficient modules it is
the zero-extension map
\((v_i)_{i\in S}\mapsto(w_i)_{i\in T}\)
with \(w_i=v_i\) for \(i\in S\) and
\(w_i=0\) otherwise. Applying the formulas
of (MLT.12), (MLT.19), and (MLT.26)
in each active position proves that all
face insertion, localization, quotient,
Tor boundary, unit, and source squares
commute exactly. A zero amplitude keeps
its declared active position under each
map. This is an explicit change of
face label; it is not an identification
of the two tagged carriers.

More generally any original cochain
arrow between source-labelled polynomial
subcomplexes, retaining their specified
source restrictions, localizes termwise
and induces the quotient arrow:
\([u^{-r}x]\mapsto[u^{-r}f(x)]\).
This is well-defined because \(f\)
preserves the original polynomial
submodules and is \(A\)-linear.
The differential and the boundary
calculation in (MLT.16) commute with
this arrow by direct substitution.
Thus proper source labels and all
their inclusion arrows remain attached.
The explicit cohomology identification
(MLT.5) is for the original full
polynomial complex; it is not
asserted for an arbitrary proper
source restriction.

\paragraph{The actual primary connection and its full saturation.}
Now take the original full primary
\(h(s)=(s-\rho)^m\), put \(N=m+1\),
and retain
\[
 c=-\frac{(-\rho)^N}{N},\qquad
 \Phi_k=\sum_i\frac{(s_i-\rho)^N}{N}+kc,\qquad
 a=kc,\qquad d=m.
                                                               \tag{MLT.33}
\]
For an actual nontrivial zero \(\rho\),
\(\rho\ne0\), so \(c\ne0\), and \(a\ne0\)
over \(\mathbb C\). The statements about
saturation below use this actual nonzero
constant, with no deletion of it.
In the original monomial basis let
\[
 B_s(s^b)=\frac{(b+1)s^b-b\rho s^{b-1}}{N}
       \quad(0\le b<m),\qquad
 B_k=\sum_{i=1}^k1^{\otimes(i-1)}
                   \otimes B_s\otimes1^{\otimes(k-i)}.
                                                               \tag{MLT.34}
\]
The second term is zero for \(b=0\).
This formula means that its coefficient
column is column \(b\) of \(B_s\).
The original monic division is
\(\Phi_h(s)s^b=h(s)(s-\rho)s^b/N+cs^b\).
Subtract its \(L\)-image to get
\([\Phi_hs^b]=[cs^b-uB_s(s^b)]\).
Therefore the original degree-one
connection \(\partial_u-\Phi_h/u^2\)
on cohomology is
\(\partial_u-c/u^2+B_s/u\).
On the two-term complex its degree-zero
operator is \(\partial_u-\Phi_h/u^2+1/u\):
\([L,\partial_u]=-\partial_s\) and
\([L,-\Phi_h/u^2]=-h/u\) give the
exact chain identity between the two
degrees. Tensoring these identities
gives the original connection on \(M\),
\[
               \nabla=\partial_u-\frac{a}{u^2}I
                                      +\frac{B_k}{u}.     \tag{MLT.35}
\]
It acts on coefficient columns.
For the ordered row \(e=(e_{\mathbf b})\)
of basis vectors, the equivalent
module convention is
\(\partial e=e(-aI/u^2+B_k/u)\).
If one writes the basis vectors as a
column instead, this matrix is transposed;
no such transpose is suppressed below.

Let \(\mathcal D_A=\mathbb C\langle u,\partial\rangle/
(\partial u-u\partial-1)\). Formula (MLT.35)
defines a left \(\mathcal D_A\)-module
structure on \(M\), because it satisfies
the Leibniz rule
\(\nabla(um)=m+u\nabla m\).
For \(r\ge0\) put \(f_{r,j}=u^{-r}e_j\),
where \(j\) runs over the original ordered
multi-indices. Direct differentiation gives
\[
 \begin{gathered}
 u\partial f_{r,j}
     =-a f_{r+1,j}
       +\sum_\ell(B_k)_{\ell j}f_{r,\ell}
       -r f_{r,j},\\
 f_{r+1,j}=-a^{-1}
       \left(u\partial f_{r,j}
        -\sum_\ell(B_k)_{\ell j}f_{r,\ell}
        +r f_{r,j}\right).
 \end{gathered}                                             \tag{MLT.36}
\]
At \(r=0\) the entire layer \(f_{0,j}=e_j\)
lies in \(L\). If the entire \(r\)-th
layer lies in \(\mathcal D_A L\), the
right side of (MLT.36) lies there for
every \(j\), since it uses its derivative,
multiplication by \(u\), and complex
linear combinations of that same layer.
Induction proves that all negative
powers of all original basis vectors
belong to \(\mathcal D_A L\). They
generate \(M\) as an \(A\)-module,
and \(M\) is already stable under
\(\mathcal D_A\). Thus the exact
saturation is
\[
                         \mathcal D_A L=M.               \tag{MLT.37}
\]
The proof uses all columns at every
step and makes no assumption that
an arbitrary intermediate submodule
is \(B_k\)-stable.

The lattice \(L\) is not stable under
this connection: \(\partial e_j\)
has coefficient \(-a e_j\ne0\) at
the pole \(u^{-2}\), which cannot
cancel the \(u^{-1}\) coefficient.
Therefore the projection \(M\to Q\)
does not define an induced quotient
\(\mathcal D_A\)-module structure.
If it did, \(e_j\mapsto0\) would
force \(\partial e_j\mapsto0\),
contradicting \(\partial e_j\notin L\).
This is the exact non-descent
statement for the original connection.
The underlying module \(Q\) can,
after specifying the same coefficient
frame, be given the different
ordinary principal-parts derivative
\(\partial_0[u^{-r}v]=-r[u^{-r-1}v]\).
That derivative descends because
ordinary differentiation preserves
\(L\). It is not substituted for
the connection (MLT.35).

\paragraph{Its exact two-pole obstruction on that same quotient.}
The difference of the original
connection and the coefficient
derivative is the \(A\)-linear
Laurent matrix \(C(u)=-aI/u^2+B_k/u\).
Its obstruction to preserving the
lattice is the actual map
\[
 \Omega:L\longrightarrow Q,\qquad
 \Omega(v)=[C(u)v].
 \quad
 \Omega(v_0+uv_1+u^2w)
   =\left[-\frac{av_0}{u^2}
            +\frac{B_kv_0-av_1}{u}\right].               \tag{MLT.38}
\]
The last equality follows by multiplying
out and discarding only the resulting
nonnegative powers, which lie in \(L\).
It proves \(\ker\Omega=u^2L\), since
the \(u^{-2}\) coefficient first forces
\(v_0=0\), and then the \(u^{-1}\)
coefficient forces \(v_1=0\).
For any class \([x/u^2+y/u]\),
the choices
\[
 v_0=-x/a,\qquad v_1=-B_kx/a^2-y/a
 \quad\text{give}\quad
 \Omega(v_0+uv_1)=[x/u^2+y/u].
 \qquad L/u^2L\xrightarrow[\Omega]{\ \sim\ }Q[u^2].
                                                               \tag{MLT.39}
\]
Every class in \(Q[u^2]\) has that
unique two-term form by (MLT.11).
This proves surjectivity and the
displayed isomorphism, including
its full \(B_k\) coefficient.
The equation \(\Omega(uv)=u\Omega(v)\)
also follows either from its
Laurent-matrix definition or from
(MLT.38), so the isomorphism is
one of \(A\)-modules.
This computes the original
non-descent obstruction through
both pole layers while (MLT.14)
retains the first derived marked
fibre.

\paragraph{Connection and periods with the same unit and lattice.}
Put \(y_i=s_i-\rho\) only as the
original displayed coordinate
translation. The Pascal matrix is
\[
 P_{ab}=
 \begin{cases}\binom ba\rho^{b-a}&a\le b,\\0&a>b,\end{cases}
 \quad
 (P^{-1})_{ab}=
 \begin{cases}\binom ba(-\rho)^{b-a}&a\le b,\\0&a>b,\end{cases}
 \quad P_k=P^{\otimes k}.                                 \tag{MLT.40}
\]
The binomial theorem in both directions
proves the inverse and the equality
of ordered rows \(e_s=e_yP_k\).
It also gives
\(P_kB_kP_k^{-1}
=\operatorname{diag}(\sum_i(b_i+1)/N)\)
by applying (MLT.34) to \(y\)-monomials.
This is a constant coefficient
isomorphism of the entire localization
sequence, with Tor transport (MLT.19).
It preserves the stated top form.

For a branch of \(u^{1/N}\) on
a simply connected sector with \(u\ne0\),
retain all the original period constants
and orientations:
\[
 \begin{gathered}
 \zeta=e^{2\pi i/N},\qquad
 \beta_b=(b+1)/N,\qquad
 d_b=N^{\beta_b-1}e^{\pi i\beta_b}\Gamma(\beta_b),\\
 W_{jb}=\zeta^{j(b+1)}-1\quad(1\le j\le m,\ 0\le b<m),
 \qquad D(u)=\operatorname{diag}(d_bu^{\beta_b}),\\
 \Pi_k(u)=e^{a/u}W^{\otimes k}D(u)^{\otimes k}P_k .
 \end{gathered}                                             \tag{MLT.41}
\]
Here \(\Gamma(\beta)=\int_0^\infty
t^{\beta-1}e^{-t}\,dt\), which is
strictly positive for \(0<\beta<1\).
The original rays have angles
\((\arg u+\pi)/N+2\pi j/N\), oriented
outwards, and contour
\(\Gamma_j=-\ell_0+\ell_j\).
The \(k\)-fold contour has their
ordered product orientation.
On the ray, substitution
\(t=r^N/(N|u|)\) computes its
\(y^bdy\) integral as
\(d_bu^{\beta_b}\zeta^{j(b+1)}\).
Subtracting the zeroth ray gives
the stated entry of \(WD\).
The translation from the original
\(s\)-contour to \(\rho+\Gamma_j\)
is justified by the entire integrand:
on a remote connector
\(s=Re^{i\vartheta}+t\rho\),
\(0\le t\le1\), the real part
of its exponent is
\(-R^N/(N|u|)+O(R^{N-1})+O(1)\),
uniformly in \(t\). Its length is
at most \(|\rho|\), so every
polynomially weighted connector
integral tends to zero. Finite
junction connectors cancel between
the two relative rays. Cauchy's
theorem and this estimate prove
the translation equality, retaining
the scalar \(e^{c/u}\). The same
decay proves absolute integrability,
so Fubini gives its ordered tensor
product, exactly (MLT.41).

The matrix \(W\) is invertible:
a relation of its columns defines
a polynomial
\(\sum_{r=1}^m c_r(z^r-1)\)
of degree at most \(m=N-1\)
vanishing on all \(N\)-th roots
of unity, including \(1\).
It is therefore zero, and its
coefficients give every \(c_r=0\).
The matrices \(D(u)\) and \(P_k\)
are invertible as well, by their
nonzero diagonal and (MLT.40).
Differentiating their explicit
entries gives the exact analytic
identity
\[
 \Pi_k'=\Pi_k(-aI/u^2+B_k/u),\qquad
 \frac{d}{du}(\Pi_k x)=\Pi_k\nabla x,\qquad
 \text{period of source coordinate }q
       =\Pi_kU_k^{-1}q.                                  \tag{MLT.42}
\]
The first equation uses
\(D_k'=D_k\operatorname{diag}
(\sum_i(b_i+1)/N)/u\) and (MLT.40);
the second is its product rule;
the third is the actual coordinate
\(q=U_kx\) in (MLT.26).
No coefficient unit or primitive
constant has been removed.

Multiplication by \(\Pi_k\) on that
sector identifies the \(A\)-module
sequence \(L\subset M\to Q\) with
\(\Pi_kL\subset\Pi_kM\to\Pi_kM/\Pi_kL\).
Injectivity follows from pointwise
invertibility and the fact that
a Laurent polynomial vector zero
on a sector is identically zero.
These are the specified
\(A\)-modules of functions;
they have not been enlarged
to all holomorphic sections.
Under this identification the
Tor boundary is still
\[
 [u^{-1}\Pi_k v]\longmapsto
       \Pi_k v\bmod u\,\Pi_k L,                           \tag{MLT.43}
\]
because multiplication by \(\Pi_k\)
commutes with \(u\). Thus the
analytic period module, the
original lattice, and the marked
Tor class belong to one exact
diagram.

Finally the full unit's connection
defect is
\[
 [\nabla,U_k]=\frac{[B_k,U_k]}{u}.
                                                               \tag{MLT.44}
\]
It follows by applying both sides
to an arbitrary coefficient column:
\(U_k\) is constant in \(u\),
and the scalar double-pole term
commutes with it.
For the primary Taylor coefficients,
put \(a_j=v_h^{(j)}(\rho)/j!
=g^{(m+j)}(\rho)/(m+j)!\) for
\(0\le j<m\); division of the
Taylor series of \(g\) by
\((s-\rho)^m\) proves this equality.
Set \(a_{\mathbf j}=\prod_i a_{j_i}\)
for \(0\le j_i<m\), and let \(J_i\)
be multiplication by \(y_i\) in
the original ordered tensor
\(y\)-basis modulo all \(y_i^m\).
Then \(U_k=\sum_{\mathbf j}a_{\mathbf j}
\prod_iR_i^{j_i}\), with
\(R_i=P_k^{-1}J_iP_k\),
the diagonal formula in (MLT.40)
gives the full expression
\[
 [\nabla,U_k]
   =\frac1{Nu}\sum_{\mathbf j}
       |\mathbf j|a_{\mathbf j}
                   \prod_iR_i^{j_i}.                    \tag{MLT.45}
\]
Indeed \(J_i\) raises the \(i\)-th
exponent by one, so the diagonal
residue difference is \(1/N\);
iteration gives \(|\mathbf j|/N\).
At a truncation boundary both
sides are zero. Formula (MLT.45)
keeps every original Taylor
coefficient and its exact
nonhorizontal term.
The results above construct the
localization triangle, its entire
marked derived fibre, and its
actual source, support, unit,
connection, and period maps.
They make no identification
with a finite-field Frobenius
action and give no RH verdict.
